
\section{Ideal gaz holat tenglamasi}

\begin{enumerate}

\item $1 \, cm^3$ hajmdagi gazning normal sharoitdagi molekulalar sonini hisoblang ($sm^{-3}$).
\begin{enumerate}[label=\Alph*)]
\item $2,1 \cdot 10^{18}$
\item $4,6 \cdot 10^{18}$
\item $1,5 \cdot 10^{19}$
\item $2,7 \cdot 10^{19}$
\end{enumerate}

\item $1 \, m^3$ kisloroddagi molekulalar soni normal sharoitda qancha bo‘lishini aniqlang.
\begin{enumerate}[label=\Alph*)]
\item $8,1 \cdot 10^{25}$
\item $2,7 \cdot 10^{25}$
\item $1,8 \cdot 10^{23}$
\item $3,8 \cdot 10^{25}$
\end{enumerate}

\item Normal sharoitda $1 \, m^3$ gazdagi molekulalar soni... bilan aniqlanadi.
\begin{enumerate}[label=\Alph*)]
\item Bolstman doimiysi
\item molyar gaz doimiysi
\item Avogadro soni
\item Loshmidt soni
\end{enumerate}

\item Hozirgi zamon texnikasi yordamida $1 \, pPa$ vakuum hosil qilish mumkin. Ana shunday $1 \, sm^3$ vakuumda 300 K temperaturada nechta gaz molekulasi qoladi?
\begin{enumerate}[label=\Alph*)]
\item $\approx 22$
\item $\approx 242$
\item $\approx 424$
\item $\approx 214$
\end{enumerate}

\item O‘lchami $2,5 \times 4 \times 3 \, m^3$ bo‘lgan idishdagi gazning bosimi 831 mm.sim.ust, temperaturasi $27^\circ C$ bo‘lsa, undagi molekulalar soni nechta?
\begin{enumerate}[label=\Alph*)]
\item $1,8 \cdot 10^{25}$
\item $6 \cdot 10^{26}$
\item $2,4 \cdot 10^{23}$
\item $8 \cdot 10^{26}$
\end{enumerate}

\item Normal sharoitdagi hajm $1 \, l$ bo‘lgan havo molekulalarining sonini aniqlang.
\begin{enumerate}[label=\Alph*)]
\item $2,7 \cdot 10^{22}$
\item $6 \cdot 10^{23}$
\item $2 \cdot 10^{20}$
\item $2,7 \cdot 10^{15}$
\end{enumerate}

\item Agar $15 \, l$ hajmli idishda $1,8 \cdot 10^{24}$ molekula bo‘lsa, 0,5 MPa bosimdagi bu gazning temperaturasi qanday (K)?
\begin{enumerate}[label=\Alph*)]
\item 123
\item 504
\item 234
\item 301
\end{enumerate}

\item Hajmi $0,8 \, m^3$ bo‘lgan gazning 300 K temperaturadagi bosimi $2,8 \cdot 10^5$ Pa. Shu gaz $3,2 \cdot 10^5$ Pa bosimda $1,4 \, m^3$ hajmni egallasa, uning temperaturasi qanday bo‘ladi (K)?
\begin{enumerate}[label=\Alph*)]
\item 150
\item 165
\item 330
\item 600
\end{enumerate}

\item Gazning $7,2 \cdot 10^5$ Pa bosim va 288 K temperaturadagi hajmi $0,6 \, m^3$ ga teng. Agar bosim $2,25 \cdot 10^5$ Pa bo‘lsa, xuddi shuncha gaz massasi qanday temperaturada $1,6 \, m^3$ hajmni egallaydi (K)?
\begin{enumerate}[label=\Alph*)]
\item 330
\item 120
\item 240
\item 480
\end{enumerate}

\item Bosimi 6 atm va temperaturasi 293 K bo‘lganda gaz 586 l hajmni egallaydi. Xuddi shu gaz massasining 248 K temperatura va 4 atm bosimda qanday hajmni egallashini toping (l).
\begin{enumerate}[label=\Alph*)]
\item 729
\item 428
\item 676
\item 744
\end{enumerate}

\item 8 litr hajmli ideal gaz 900 K haroratda 1,5 atm bosim hosil qiladi. Shu gazning bosimi 4 litr hajmda 1 atm bo‘lishi uchun uning harorati qanday bo‘lishi kerak ($^\circ C$)?
\begin{enumerate}[label=\Alph*)]
\item 30
\item 27
\item 127
\item 237
\end{enumerate}

\item Suv havzasi tubidan sirtiga ko‘tarilgan havo pufagining diametri 1,5 marta ortdi. Suv tubidagi harorat $5^\circ C$ sirtidagi esa $27^\circ C$ bo‘lsa, suv havzasining chuqurligi necha metr?
\begin{enumerate}[label=\Alph*)]
\item 36
\item 25
\item 21,3
\item 60
\end{enumerate}

\item Hajm va absolyut harorati 2 martadan oshirilsa, ideal gazning bosimi qanday o‘zgaradi?
\begin{enumerate}[label=\Alph*)]
\item o‘zgarmaydi
\item 2 marta ortadi
\item 4 marta kamayadi
\item 4 marta ortadi
\end{enumerate}

\item Agar gaz temperaturasi 300 K ga ko‘tarilganda, bosim va hajm 2 martadan ortgan bo‘lsa, gazning oxirgi temperaturasi necha kelvinga teng?
\begin{enumerate}[label=\Alph*)]
\item 350
\item 400
\item 600
\item 900
\end{enumerate}

\item Shar zond $27^\circ C$ temperaturada 105 kPa bosimgacha gaz bilan to‘ldirilgan. Shar 80 kPa bosimli balandlikka ko‘tarilgandan so‘ng hajmi 5\% ortgan bo‘lsa, u yerdagi haroratni (K) toping.
\begin{enumerate}[label=\Alph*)]
\item 125
\item 445
\item 500
\item 255
\end{enumerate}

\item Ideal gazning absolyut temperaturasi 2 marta ortganda uning bosimi 25\% ortdi. Bunda hajmi necha marta o‘zgaradi?
\begin{enumerate}[label=\Alph*)]
\item 1,6 marta ortadi
\item 1,6 marta kamayadi
\item 0,25 marta ortadi
\item 0,25 marta kamayadi
\end{enumerate}

\item Rezina qayiqqa $7^\circ C$ temperaturada 108 kPa ishchi bosimgacha dam berildi. Harorat $37^\circ C$ ga yetganda qayiqning yorilish xavfi bormi?
\begin{enumerate}[label=\Alph*)]
\item Bor
\item Yo‘q
\item Qayiqning massasiga bog‘liq
\item Ma'lumot yetarli emas
\end{enumerate}

\item Radiusi 2 sm bo‘lgan yupqa devorli rezina sharga 0,1 MPa bosimda va $20^\circ C$ temperaturada havo qamalgan. Uni $4^\circ C$ temperaturali ko‘lga 20 m chuqurlikkacha tushirilsa, shar radiusi qancha bo‘ladi (sm)?
\begin{enumerate}[label=\Alph*)]
\item 0,75
\item 2,12
\item 1,37
\item 1,78
\end{enumerate}

\item Ideal gazning harorati 4 marta ortganda uning hajmi 2 marta ortsa, bosimi qanday o‘zgaradi?
\begin{enumerate}[label=\Alph*)]
\item 4 marta ortadi
\item o‘zgarmaydi
\item 2 marta kamayadi
\item 2 marta ortadi
\end{enumerate}

\item Ideal gazning mutlaq harorati 3 marta oshganda bosimi 1,5 marta oshgan bo‘lsa, uning zichligi qanday o‘zgargan?
\begin{enumerate}[label=\Alph*)]
\item 1,5 marta oshgan
\item 1,5 marta kamaygan
\item o‘zgarmagan
\item 2 marta kamaygan
\end{enumerate}

\end{enumerate}



\begin{table}[h!]
\centering
\begin{tabular}{|c|c|c|c|c|c|c|c|c|c|}
\hline
№ & J. & № & J. & № & J. & № & J. & № & J. \\ \hline
1 & D & 5 & D & 9 & C & 13 & A & 17 & A \\ \hline
2 & B & 6 & A & 10 & D & 14 & B & 18 & C \\ \hline
3 & D & 7 & D & 11 & B & 15 & D & 19 & D \\ \hline
4 & B & 8 & D & 12 & C & 16 & A & 20 & D \\ \hline
\end{tabular}
\end{table}